\subsection{The last permanent seven-plane: a second prolongation}
\begin{finding}{P-19}{the final permanent seven-plane has border rank at least fourteen}
Over every algebraically closed field of characteristic zero, the
permanent residual tensor for the rectangle subset
\[
 S=\{0,1,2,6,7,18,19\}
\]
has border rank at least $14$.
\end{finding}
\begin{proof}
We give the additional closed necessary condition, the complete reduction
to eight fixed extensions, and their integer certificates.
Use $A=k^7$ and $B=C=k^{11}$, with respective bases
$a_0,\ldots,a_6$ and $x_0,\ldots,x_{10}$.  The variables are the entries
of a three-by-four rectangle, numbered in row order, with $x_{11}$
omitted.  For compactness write $xy$ for $x\otimes y+y\otimes x$ in
$B\otimes C$, throughout this proof.  The tensor is
\[
 T=\sum_{q=0}^6a_q\otimes Q_q,
 \qquad
 \begin{aligned}
 Q_0&=x_0x_5+x_1x_4,& Q_1&=x_0x_6+x_2x_4,\\
 Q_2&=x_0x_7+x_3x_4,& Q_3&=x_0x_9+x_1x_8,\\
 Q_4&=x_0x_{10}+x_2x_8,& Q_5&=x_4x_9+x_5x_8,\\
 Q_6&=x_4x_{10}+x_6x_8.
 \end{aligned}
\]
Thus the displayed notation does not presume a decomposition with equal
second and third factors.  Let $U=T(C^*)\subset A\otimes B$;
its eleven contraction vectors are independent.

\paragraph{A closed higher-prolongation necessary condition.}
For a subspace $E\subset A\otimes B$, use the symmetric-tensor
subspaces to define
\begin{align*}
 \mathcal P_{21}(E)&=(A\otimes E)\cap(\operatorname{Sym}^2A\otimes B),\\
 \mathcal P_{12}(E)&=(E\otimes B)\cap(A\otimes\operatorname{Sym}^2B),\\
 \mathcal P_{22}(E)&=(A\otimes E\otimes B)
               \cap(\operatorname{Sym}^2A\otimes\operatorname{Sym}^2B).
\end{align*}
Factor permutations place these intersections in $A^{\otimes2}\otimes B$,
$A\otimes B^{\otimes2}$, and $A^{\otimes2}\otimes B^{\otimes2}$,
respectively.  In the last expression $E$ occupies the second $A$ and
first $B$ positions.
For each fixed dimension $r$ of $E$, each condition
$\dim\mathcal P_{ij}(E)\ge r$ is closed on
$\operatorname{Gr}(r,A\otimes B)$.  Indeed it is a kernel-dimension
condition for the map from the corresponding tensor product of the
rank-$r$ tautological bundle to the fixed ambient tensor space modulo
the displayed symmetric subspace.  Both source and target have fixed
ranks, so it is determinantal.

If $E$ is spanned by $r$ independent pure vectors
$u_\ell\otimes v_\ell$, its prolongations contain, respectively,
\[
 u_\ell^{\otimes2}\otimes v_\ell,\qquad
 u_\ell\otimes v_\ell^{\otimes2},\qquad
 u_\ell^{\otimes2}\otimes v_\ell^{\otimes2}
 \quad(1\le\ell\le r).
\]
Each of these $r$-tuples is independent.  Choose linear forms nonzero
on all the finitely many $u_\ell$ and $v_\ell$, and contract the added
factors; the result is the original independent tuple multiplied by
nonzero scalars.  Such forms exist because $k$ is infinite.
The open condition that the $r$ pure $AB$ vectors be independent is
nonempty and dense in the parameter space of $r$ pure $ABC$ summands
when $r\le\dim(A\otimes B)$.
Consequently, impose $T(C^*)\subset E$ and all three closed
prolongation conditions above in the product with the projective
Grassmannian.  Its projection is closed and contains the dense
parametrized rank-$r$ set.  This is the same proper-incidence argument
as in R04, now with the additional $\mathcal P_{22}$ condition.
It proves the needed implication
\begin{equation}
 \underline R(T)\le13\ \Longrightarrow\
 \exists E\supset U,\quad \dim E=13,\quad
 \dim\mathcal P_{21}(E),\dim\mathcal P_{12}(E),
       \dim\mathcal P_{22}(E)\ge13.
 \label{eq:P19-necessary}
\end{equation}
In particular, the third condition is imposed directly on the fixed
ambient intersection.  No assertion of continuity for a varying
first-prolongation kernel is used.

\paragraph{The complete first-stage enumeration.}
Apply the row/column split torus with negative quadratic weights on
$A$ and variable weights on $B,C$, as in P-15.  The closed incidence
in \eqref{eq:P19-necessary} is torus invariant.  A separating
cocharacter limit therefore makes $E$ torus invariant without losing
any of its conditions.  The quotient $(A\otimes B)/U$ has $47$
nonzero weight blocks and dimension $66$.  Fourteen blocks have
capacity at least two.  Since $\dim(E/U)=2$, there are
\[
 14+\binom{47}{2}=1095
\]
possible supports, consisting of a two-plane in one weight or lines
in two different weights.

The checker \texttt{check\_per7\_last\_certificate.py} constructs these
blocks in lexicographic $(q,v)$ order.  As in P-15, in every image-weight
block it omits the last coordinate, whose image coefficient is nonzero,
to specify a basis of the quotient.  The baseline wedge matrices are
\[
\begin{array}{c|c|c|c}
 &\text{matrix on }U&\text{rank}&\text{maximum allowed increment}\\\hline
 \Phi_A&231\times77&70&8\\
 \Phi_B&385\times121&115&15.
\end{array}
\]
The allowed total ranks are $78$ and $130$, respectively.
Saturating each chosen support by its entire weight blocks, a
prolongation kernel of dimension at most twelve excludes $971$ of the
$1095$ supports.  Every surviving support consists of two distinct
weight lines; there are $124$ such pairs.

For each line, the projected columns in one orientation have pairwise
distinct output weights.  The complete linear-arrangement algorithm
proved in P-15 gives the uniform minimum number of nonzero columns.
For a pair, the sum of those minima minus the number of common output
weights is a lower bound for the rank increment.  This excludes another
$116$ pairs.  The eight remaining pairs, in the fully specified block
numbering of the checker, are
\[
 (1,23),(5,23),(19,23),(20,23),(21,23),(22,23),(23,24),(23,25).
\]
All the saturation ranks, quotient annihilators, and line arrangements
in these two stages are recomputed over $\mathbb Q$ by the attached
self-contained checker; no earlier search file or numerical rank is
used.

\paragraph{Joint refinement gives eight fixed extensions.}
Block $23$ is one dimensional and has the coordinate representative
$a_2\otimes x_8$.  Thus every remaining extension contains
\[
 U_0=U+\langle a_2\otimes x_8\rangle,\qquad \dim U_0=12.
\]
Its wedge-map ranks are $72$ and $123$.  The remaining line may
therefore add at most six $A$-columns and seven $B$-columns.
Quotient by the actual images on $U_0$, and apply the joint
$7+11$-column arrangement enumeration.  The required zero counts are
at least one $A$-column and four $B$-columns.

Only blocks $1$ and $5$ have parameters; their quotient dimensions
are four, with the following ordered coordinate representatives:
\[
\begin{array}{c|l|c}
 \text{block}&\text{representatives}&\text{only surviving parameter line}\\\hline
 1&a_0x_1,\ a_1x_2,\ a_2x_3,\ a_5x_9&(0,0,1,0)\\
 5&a_0x_5,\ a_1x_6,\ a_2x_7,\ a_3x_9&(0,0,1,0).
\end{array}
\]
Here $a_qx_v$ means $a_q\otimes x_v$.  In each block the exact joint
linear arrangement has six strata, of which the displayed line alone
satisfies both rank conditions.  This uses the entire projective
parameter space; the other five strata have not been replaced by
sampled parameter values.  The six other blocks are already one
dimensional, with representatives
\[
 a_2x_1,\ a_2x_2,\ a_2x_5,\ a_2x_6,\ a_2x_9,\ a_2x_{10}
 \quad\text{for blocks }19,20,21,22,24,25.
\]
It follows that every torus-invariant $E$ satisfying the two first
prolongation conditions is one of the following eight fixed spaces:
\begin{equation}
 E_v=U+\langle a_2\otimes x_8,\ a_2\otimes x_v\rangle,\qquad
 v\in\{1,2,3,5,6,7,9,10\}.
 \label{eq:P19-eight}
\end{equation}
The joint linear equations and their unique surviving parameter
vectors are included explicitly in \texttt{per7\_last\_certificate.json}.

\paragraph{The eight degree $(2,2)$ obstructions.}
For each space in \eqref{eq:P19-eight}, the map
\[
 \Phi_B:E_v\otimes B\longrightarrow A\otimes\Lambda^2B
\]
is a $385\times143$ integer matrix $M_{12,v}$.
The certificate provides a $130\times130$ unit minor and an integer
matrix $N_v$ of size $143\times13$ satisfying
\[
 M_{12,v}N_v=0,
\]
together with a $13\times13$ unit minor of $N_v$.
Thus $\operatorname{rank}M_{12,v}=130$, and the columns of $N_v$ form
an exact basis of its kernel $\mathcal P_{12}(E_v)$.

We use the symmetric-tensor basis
$x_i\otimes x_i$ and
$x_i\otimes x_j+x_j\otimes x_i$ for $i<j$.
The checker embeds each kernel vector in $A\otimes B\otimes B$,
verifies symmetry under exchanging its two $B$ indices, and reads its coefficients
in this basis.  There are $\binom{12}{2}=66$ symmetric coordinates.
The resulting thirteen matrices in $A\otimes\operatorname{Sym}^2B$
give the next wedge map
\[
 M_{22,v}:A\otimes\mathcal P_{12}(E_v)
       \longrightarrow\Lambda^2A\otimes\operatorname{Sym}^2B,
 \qquad \text{matrix size }1386\times91.
\]
Its kernel is exactly $\mathcal P_{22}(E_v)$: first intersect
$A\otimes E_v\otimes B$ with $A^{\otimes2}\otimes\operatorname{Sym}^2B$,
then impose symmetry in the two $A$ factors.  The two successive
intersections equal the defining simultaneous intersection.

Each $M_{22,v}$ has a certified $79\times79$ unit minor.
For clarity, all twenty-four final integer determinants are listed here,
with the selected minor order in square brackets:
\[
\begin{array}{c|r|r|r}
 v&\det(M_{12,v}[130])&\det(N_v[13])&\det(M_{22,v}[79])\\\hline
 1&1&1&1\\
 2&-1&1&1\\
 3&1&1&1\\
 5&-1&1&-1\\
 6&1&1&1\\
 7&-1&1&1\\
 9&1&1&1\\
 10&-1&1&1
\end{array}.
\]
The certificate contains every selected row and column and every entry
of $N_v$.  The checker evaluates the determinants over $\mathbb Z$;
its auxiliary modular elimination only selects the minors.
Consequently
\[
 \dim\mathcal P_{22}(E_v)
   \le91-79=12<13
\]
for each of the eight spaces.  This contradicts
\eqref{eq:P19-necessary} in every possible case and proves the finding.
\end{proof}

\paragraph{Scope and dictionary.}
The last twenty-four minors are units, but the preceding finite
arrangement reduction was checked over $\mathbb Q$.  P-19 therefore
claims characteristic zero, without an unverified extension to every
prime.  Geometrically it uses determinantal closed subschemes of a
relative Grassmannian, its tautological bundle, and a proper incidence
projection, with a split-torus specialization.  The extra condition is
the next symmetric prolongation of that bundle; it is a kernel rank
condition in a fixed ambient locally free sheaf.  Exterior and symmetric
powers supply the standard $\lambda$-ring and Schur-functor dictionary
at the common SGA/FGA locators of the paper.  No purity, $\ell$-adic
weight, or motivic claim is needed to transfer this proof.
